\section{Evidence Provenance and Scope}
\label{app:evidence-provenance}

\subsection{FULL task portfolio}

The task summary and condition-level records are archived in the paper-results package dated 2026-11-13. The summary reports eleven task cases, 32 required/control records, 30 passed required conditions, and two non-passing T10 frozen controls. The result set is a snapshot of completed studies, not a repeated-seed statistical evaluation. The original task-specific evaluator and source trace remain the authoritative references for every row.

\subsection{Nine-segment PHS-2D physical case}

The principal PHS-2D mechanism-design case uses the nine-segment, four-tendon Round-3 evidence package. The final structure was checked across ten loading conditions, with independent-process reproduction, half-time-step evaluation, and initial-state perturbation checks. The reported summary values are a minimum loaded-point displacement of 11.155 mm, a maximum projected tip displacement of 3.068 mm, and a maximum displacement ratio of 0.1093. These values must remain attached to the nine-segment evidence package; older seven-segment snapshots are not combined with them.

\subsection{PHS-2D local sensitivity audit}

The gradient audit is a separate fixed-topology local sensitivity record for a nine-segment, four-tendon passive model. It uses one common Young's-modulus direction applied across the nine segments, with normalized perturbation $h=0.001$. Six loading branches and three task metrics (loaded-point displacement, horizontal tip displacement, and their ratio) yield 18 task-metric derivative pairs. The maximum relative discrepancy between the KKT-based implicit task-metric derivatives and finite differences obtained by re-equilibrating the perturbed designs is 0.1191223\%.

The 0.1191223\% figure applies to those 18 task-metric derivative pairs only. It is not the maximum error over every internal generalized-coordinate derivative, not 18 independent design-parameter gradients, and not a global error bound. The implicit calculation uses numerically estimated partial derivatives in the KKT system; the finite-difference values are derived from actual saved positive/negative perturbed equilibria. The audit does not establish a universal descent guarantee or high-fidelity closed-loop improvement.

\subsection{Historical C and new A/B provenance}

C is the historical FULL record. A and B are new runs. The manuscript retains the historical configuration and source identifiers for C and does not label the comparison as a fully controlled randomized experiment when model, context, budget, or protocol versions differ. The first batch is four tasks and eight new A/B runs. Results are inserted only after the corresponding runs and configuration review are complete.

\subsection{Physical-robot demonstrations}

Real-world robot results are maintained as a separate evidence source. They are not inferred from MuJoCo trajectories and are not included in the simulation FULL counts. The final physical-validation section will cite the existing hardware records and selected demonstration videos/frames.
